\chapter{Numerical Simulation of the Atom-Photon Gate}\label{appendix:qutip_simulation}

This appendix describes the time-domain master-equation simulations used to determine the optimal Brewster window attenuation (section~\ref{appendix:bw_scan}) and to characterize the gate fidelity as a function of mean photon number (section~\ref{appendix:photon_number_scan}). All simulations were implemented in Python using the QuTiP framework \cite{qutip5}.

\section{Simulation Model}\label{appendix:sim_model}

The simulation models the full open quantum system consisting of a single $^{87}$Rb atom coupled to a birefringent Fabry-P\'{e}rot cavity supporting two polarization modes ($\pi$ and $V$). The atom is described by a four-level Hilbert space spanned by $\{\ket{0}, \ket{1}, \ket{\text{dark}}, \ket{e}\}$, where $\ket{0}$ and $\ket{1}$ are the clock qubit states, $\ket{\text{dark}}$ captures off-resonant decay channels, and $\ket{e}$ is the excited state coupled to $\ket{1}$ via the cavity mode. Each cavity mode is truncated to a Fock space of dimension $N_{\text{ph}}$, yielding a total Hilbert space of dimension $N_{\text{ph}} \times N_{\text{ph}} \times 4$.

\subsection{Hamiltonian}

In the rotating frame, the system Hamiltonian reads
\begin{align}
	\hat{H}_0 & = \Delta_{c,\pi}\, \hat{a}_\pi^\dagger \hat{a}_\pi + (\Delta_{c,\pi} + \Delta_{\text{biref}})\, \hat{a}_V^\dagger \hat{a}_V + \Delta_a\, \hat{\sigma}^\dagger \hat{\sigma},
\end{align}
with the atom--cavity interaction
\begin{align}
	\hat{H}_{\text{int}} & = g\bigl(\hat{a}_\pi \hat{\sigma}^\dagger + \hat{a}_\pi^\dagger \hat{\sigma} + \hat{a}_V \hat{\sigma}^\dagger + \hat{a}_V^\dagger \hat{\sigma}\bigr),
\end{align}
where $g$ is the vacuum coupling rate (identical for both modes), $\Delta_{c,\pi}$ is the cavity--laser detuning for the $\pi$ mode, $\Delta_{\text{biref}} = 2\pi \times \SI{500}{\mega\hertz}$ is the birefringent splitting, and $\Delta_a$ is the atom--laser detuning.

The external drive is included via input--output theory. For the CNOT basis states $\ket{\pm} = (\ket{V} \pm \ket{\pi})/\sqrt{2}$, both cavity modes are driven simultaneously with the appropriate relative phase and amplitude. The time-dependent drive Hamiltonian takes the form
\begin{align}
	\hat{H}_{\text{drive}}^{(\pm)}(t) & = \frac{i\, \mu_{fc} \sqrt{\kappa_{OC}}}{\sqrt{2}}\, \varepsilon(t) \Bigl[\bigl(\hat{a}_\pi - \hat{a}_\pi^\dagger\bigr) \pm \sqrt{T_V}\bigl(\hat{a}_V - \hat{a}_V^\dagger\bigr)\Bigr],
\end{align}
where $\varepsilon(t)$ is the temporal envelope of the input pulse, $\mu_{fc}$ is the fiber--cavity mode-matching parameter, $\kappa_{OC}$ is the outcoupling rate, and $T_V$ is the V-mode power transmission set by the Brewster window array.

\subsection{Initial States and Imperfect State Preparation}

The initial state of the system is a product of the atomic state and the cavity vacuum. For the atomic qubit state $\ket{0}$, the atom is initialized in a pure state $\hat{\rho}_0 = \ket{0}\bra{0}$. To account for imperfect optical pumping and microwave transfer, the nominal $\ket{1}$ initial state is modeled as a density matrix mixture
\begin{align}
	\hat{\rho}_1 = p_1 \ket{1}\bra{1} + (1 - p_1)\ket{0}\bra{0}, \label{eq:imperfect_prep}
\end{align}
with $p_1 = 0.88$. The full initial state for each simulation run is then $\hat{\rho}_{\text{init}} = \hat{\rho}_{\text{atom}} \otimes \ket{0_\pi}\bra{0_\pi} \otimes \ket{0_V}\bra{0_V}$.

\subsection{Dissipation}

The open-system dynamics are governed by the Lindblad master equation
\begin{align}
	\dot{\hat{\rho}} = -i[\hat{H}, \hat{\rho}] + \sum_k \mathcal{D}[\hat{L}_k]\hat{\rho},
\end{align}
with collapse operators accounting for:
\begin{itemize}
	\item Photon loss through the outcoupling mirror: $\hat{L}_{1,2} = \sqrt{\kappa_{OC}}\, \hat{a}_{\pi,V}$,
	\item Internal cavity losses (absorption, HR mirror leakage): $\hat{L}_{3,4} = \sqrt{\kappa_{\text{int}}}\, \hat{a}_{\pi,V}$, where $\kappa_{\text{int}} = \kappa - \kappa_{OC}$,
	\item Spontaneous emission back to the coupled state: $\hat{L}_5 = \sqrt{p_{e \to 1}\, \gamma}\, \ket{1}\bra{e}$,
	\item Spontaneous emission to off-resonant states: $\hat{L}_6 = \sqrt{p_{e \to \text{dark}}\, \gamma}\, \ket{\text{dark}}\bra{e}$,
\end{itemize}
with branching ratios $p_{e \to 1} = 1/15$ and $p_{e \to \text{dark}} = 14/15$ for the $F' = 1$ excited state of the D$_2$ line.

\subsection{Output Fields and Detection Probabilities}

The quantum output photon number in each polarization mode is computed via the input--output relation. The mean output field amplitude is
\begin{align}
	\alpha_{\text{out},\alpha}(t) = \mu_{fr}\, \alpha_{\text{in},\alpha}(t) + \mu_{fc}\sqrt{\kappa_{OC}}\, \langle \hat{a}_\alpha(t) \rangle,
\end{align}
where $\mu_{fr}$ is the free-space mode-matching parameter. The quantum output photon number additionally includes the contribution from intracavity fluctuations:
\begin{align}
	n_{\text{out},\alpha}^{(q)}(t) = |\alpha_{\text{out},\alpha}(t)|^2 + |\mu_{fc}|^2 \kappa_{OC}\bigl(\langle \hat{a}_\alpha^\dagger \hat{a}_\alpha \rangle - |\langle \hat{a}_\alpha \rangle|^2\bigr). \label{eq:quantum_output}
\end{align}
For the CNOT basis, the output fields are transformed via $n_{\text{out},\pm}^{(q)} = (n_{\text{out},V}^{(q)} \pm n_{\text{out},\pi}^{(q)})/2$.

The detection probabilities at the two SNSPD detectors include corrections for dark counts and finite detection efficiency $\eta$.

\subsection{CNOT Truth Table Construction}\label{appendix:truth_table}

The CNOT truth table is a $4 \times 4$ matrix whose columns correspond to the four input configurations $\{\ket{1,+},\, \ket{1,-},\, \ket{0,+},\, \ket{0,-}\}$ and whose rows correspond to the four output configurations $\{\ket{1,+},\, \ket{1,-},\, \ket{0,+},\, \ket{0,-}\}$, where the first label denotes the atomic state and the second label the photonic polarization. Each entry is the joint probability of observing a given photonic detection outcome \emph{and} the atom being found in a given qubit state after the gate pulse. Concretely, the entry for output $\ket{a', p'}$ given input $\ket{a, p}$ reads
\begin{align}
	\text{CNOT}_{(a',p'),(a,p)} = P_{\text{atom}}(a' \,|\, a,p)\;\cdot\; P_{\text{det}}(p' \,|\, a,p), \label{eq:cnot_entry}
\end{align}
where $P_{\text{atom}}(a' \,|\, a,p)$ is the probability that the atom occupies qubit state $\ket{a'}$ at the end of the pulse (extracted from the final density matrix), and $P_{\text{det}}(p' \,|\, a,p)$ is the photonic detection probability in output port $p'$. Each column is then normalized to unit probability.

This construction ensures that spontaneous scattering events---which can flip the atomic qubit state during the gate pulse---are reflected as errors in the truth table, rather than being accounted for as a separate correction factor.

The average probability for the atom to remain in its initial qubit state is
\begin{align}
	p_{\text{atom}} = \frac{1}{4}\sum_{a \in \{0,1\},\; p \in \{+,-\}} P_{\text{atom}}(a \,|\, a,p),
\end{align}
which provides a direct measure of the atomic state preservation across all input configurations.

\subsection{Simulation Parameters}\label{appendix:sim_params}

The parameters used in the simulations are summarized in Table~\ref{tab:sim_params}. These correspond to the experimentally characterized values presented in section~\ref{section:experimental_parameters_for_gate}.

\begin{table}[ht]
	\centering
	\begin{tabular}{l c c}
		\hline\hline
		Parameter                   & Symbol                  & Value                                 \\
		\hline
		Atom--cavity coupling       & $g$                     & $2\pi \times \SI{25.9}{\mega\hertz}$  \\
		Total cavity decay rate     & $\kappa$                & $2\pi \times \SI{58}{\mega\hertz}$    \\
		Outcoupling ratio           & $\kappa_{OC}/\kappa$    & $0.85$                                \\
		Atomic decay rate           & $\gamma$                & $2\pi \times \SI{6.065}{\mega\hertz}$ \\
		Birefringent splitting      & $\Delta_{\text{biref}}$ & $2\pi \times \SI{500}{\mega\hertz}$   \\
		Fiber--cavity mode-matching & $\mu_{fc}$              & $0.873$                               \\
		Free-space mode-matching    & $\mu_{fr}$              & $0.978$                               \\
		State preparation fidelity  & $p_1$                   & $0.88$                                \\
		Detection efficiency        & $\eta$                  & $0.742$                               \\
		Dark count rate             & $r_{\text{dc}}$         & $5 \times 10^{-5}$                    \\
		\hline
	\end{tabular}
	\caption{Parameters used in the master-equation simulations. The detection efficiency is the product of fiber coupling ($0.9$), SNSPD efficiency ($0.85$), and optical path transmission ($0.97$).}
	\label{tab:sim_params}
\end{table}


\section{Brewster Window Transmission Scan}\label{appendix:bw_scan}

To determine the optimal number of Brewster windows for the V-mode attenuation (cf.\ section~\ref{section:polarization_dependent_losses}), the gate fidelity was computed as a function of the V-mode transmission $T_V \in [0, 1]$.

For each value of $T_V$, the simulation was run in the CNOT basis: weak coherent pulses with $\ket{+}$ and $\ket{-}$ input polarizations were reflected from the atom-cavity system for both atomic qubit states $\ket{0}$ and $\ket{1}$. In the simulation, the V-mode transmission enters as a prefactor $\sqrt{T_V}$ on the V-mode drive amplitude, effectively modeling the attenuation introduced by the Brewster window array prior to the cavity. The Fock space truncation was set to $N_{\text{ph}} = 3$ for this scan.

From the four resulting output fields, the $4 \times 4$ CNOT truth table was constructed as described in section~\ref{appendix:truth_table}, including the atomic state probability weighting. The columns of the truth table were then normalized to unit probability. The overlap was evaluated as the column-averaged overlap between the simulated and ideal CNOT truth tables:
\begin{align}
	\text{Overlap} = \frac{1}{4} \sum_{j=1}^{4} \left| \langle U_j | V_j \rangle \right|^2, \label{eq:overlap}
\end{align}
where $\ket{U_j}$ and $\ket{V_j}$ are the $j$-th columns of the ideal and simulated truth tables, respectively.

\begin{figure}
	\centering
	\includesvg[width=\linewidth]{Figures/appendix/v_transmission_cnot_overlap.svg}
	\caption{Simulated CNOT gate overlap as a function of the V-mode cavity transmission $T_\mathrm{V}$, obtained from a master-equation simulation of the atom--photon gate protocol. The dashed line marks the optimal transmission $T_\mathrm{V}^* \approx 0.24$.}
	\label{fig:bw_scan_sim}
\end{figure}

The scan shown in Figure \ref{fig:bw_scan_sim} reveals a clear optimum at $T_\mathrm{V} \approx 0.24$, corresponding to a V-mode power attenuation of approximately \SI{76}{\percent}. Given the measured single-window V-mode transmission of $(\num{77.40} \pm \num{0.13})\,\si{\percent}$ (section~\ref{sec:brewster_windows}), the na\"ive estimate $0.7738^n$ predicts an optimum between 5 and 6 Brewster windows ($0.7740^5 \approx \num{0.277}$ and $0.7740^6 \approx \num{0.215}$). However, the directly measured multi-window transmissions of $(\num{26.80} \pm \num{0.06})\,\si{\percent}$ for 6 windows and $(\num{22.47} \pm \num{0.04})\,\si{\percent}$ for 7 windows bracket $T_\mathrm{V}^*$, indicating that 6 to 7 Brewster windows realise the optimal V-mode attenuation in practice.

\section{Mean Photon Number Scan}\label{appendix:photon_number_scan}

While the atom-photon gate is designed to operate in the single-photon regime, the use of weak coherent pulses introduces a non-zero probability of multi-photon events. This section characterizes how the gate overlap depends on the mean photon number $\bar{n}$ of the input pulses.

The mean photon number was scanned over the range $\bar{n} \in [10^{-5},\, 10]$ on a logarithmic scale. For each value of $\bar{n}$, the coherent pulse amplitude was computed from the pulse normalization as
\begin{align}
	\alpha = \sqrt{\frac{\bar{n}}{\int_0^T |\varepsilon_{\text{ref}}(t)|^2\, dt}},
\end{align}
where $\varepsilon_{\text{ref}}(t)$ is the unit-amplitude pulse envelope. The CNOT truth table was then constructed as described in section~\ref{appendix:truth_table}, including the atomic state probability weighting, and the overlap with the ideal CNOT truth table was evaluated using \eqref{eq:overlap}. The V-mode transmission was fixed at $T_V = 0.2247$ and the Fock space truncation was set to $N_{\text{ph}} = 3$ for this scan.

\begin{figure}
	\centering
	\includesvg[width=\linewidth]{Figures/appendix/photon_number_scan.svg}
	\caption{Simulated CNOT gate overlap as a function of the mean input photon number $\bar{n}$. The V-mode transmission is fixed at $T_V = 0.2247$ and the Fock space is truncated at $N_{\text{ph}} = 3$. The dashed line marks the optimal mean photon number $\bar{n}^* \approx 0.10$, at which the overlap reaches a maximum of $\approx 0.865$.}
	\label{fig:photon_number_scan}
\end{figure}

The result is shown in Figure~\ref{fig:photon_number_scan}. At very low photon numbers ($\bar{n} \lesssim 10^{-3}$), dark counts dominate the detected signal and the overlap approaches $0.5$, corresponding to random guessing. As $\bar{n}$ increases, the signal-to-noise ratio improves and the overlap rises towards the single-photon limit. For $\bar{n} \gtrsim 1$, multi-photon events begin to degrade the overlap, as the master-equation simulation captures the nonlinear response of the atom-cavity system to multi-photon input states.
%  
% It should be noted that the simulation assumes a fixed atom--cavity coupling $g$ throughout the gate pulse. In practice, spontaneous scattering events impart momentum kicks to the atom, displacing it within the intracavity standing wave and thereby reducing the effective coupling strength. This effect becomes significant at higher mean photon numbers where the scattering probability per pulse is no longer negligible, and is expected to cause the simulation to overestimate the gate overlap in this regime.

%
% \chapter{Analytical Derivations}
%
% \section{Output Light Analysis}
%
% \subsection{Problem}
% From our QuTip simulation we're able to calculate the output field $a_{out}(t)$ via input output theory:
% \begin{equation}
% 	a_{out}(t) = \mu_{fr}a_{in}(t) + \mu_{fc}\sqrt{\kappa_{oc}}\langle a\rangle(t)
% 	\label{eq:io-theory}
% \end{equation}
%
% Where all the fields mentioned in \eqref{eq:io-theory}, can be in one of the following bases: $\{|\pi\rangle,|V\rangle,|+\rangle,|-\rangle\}$.
%
% Now the question becomes, how is one able to analyze this output field and generate a truth table from this. For that we're going to calculate the mean photon number in the output field via:
% $$\bar{n}_{out;\{|\pi\rangle,|V\rangle,|+\rangle,|-\rangle\} } = \int |a_{out;\{ |\pi\rangle,|V\rangle,|+\rangle,|-\rangle  \}}|^2 dt$$
%
% For the problem we're going to be looking at, we're only going to consider the $|+\rangle$ and the $|-\rangle$ modes of our light, as these are the polarizations which are set in our CNOT basis, for which we want to calculate the fidelity for.
%
% \subsection{Analysis via detector}
% Now after having extracted the mean photon number in our output field for our two polarizations we need to consider the framework which we're using to analyze this output field. In our case that is the case of 2 independent detectors, where one detects the $|+\rangle$ polarized light and the other detects the $|-\rangle$ polarized light. Both detectors share the same efficiency $\eta$, and also the same dark count rate $p_{dc}$, which is the rate at with a single detector will fire even upon receiving no signal. The efficiency can be viewed as a beam splitter. Meaning, that if we were to have $\bar{n}$ photon in from of our detector that has an efficiency $\eta$, this can be treated as a perfect detector with $\eta=1$ but with a beam splitter in front of our perfect detector that will "take" away $(1-\eta)\bar{n}$ photons from our initial pulse. The rest of the icoming pulse (now: $\eta\bar{n}$) will however get detected perfectly. Since we're only interested in the cases where we have 1 detector event ($k=1$) upon receiving $\lambda=\eta\bar{n}$ photons, we can treat this problem via poisson statistics with:
% $$ P(k,\lambda)=\frac{\lambda^ke^{-\lambda}}{k!} = \lambda e^{-\lambda}$$
% This however doesn't take into consideration our dark counts yet.
%
% \subsection{Detector statistics}
% If our detector fires, and gives us 1 detector click, it can come from 2 sources:
% \begin{itemize}
% 	\item a dark count
% 	\item a photon
% \end{itemize}
% This we can then group up into 4 cases for each of the 2 possible sources of a click
% \begin{itemize}
% 	\item $P(\text{0 Photons, 0 dark counts}) = e^{-\lambda}\cdot e^{-p_{dc}} = e^{-(\lambda+p_{dc})}$
% 	\item $P(\text{1 Photon, 0 dark counts}) = \lambda e^{-\lambda}\cdot e^{-p_{dc}} = \lambda e^{-(\lambda+p_{dc})}$
% 	\item $P(\text{0 Photons, 1 dark count}) = e^{-\lambda}\cdot p_{dc}e^{-p_{dc}} = p_{dc}e^{-(\lambda+p_{pc})}$
% 	\item $P(\text{1 Photons, 1 dark counts}) =$ Can be ignored as this would produce 2 clicks, which we're post selecting away
% \end{itemize}
%
% Meaning that we have 2 out of our 4 possible events which then produce only 1 click, giving us the probability that each of our detectors produces 1 click:
% \begin{equation}
% 	P_{\pm, click} = (\lambda_\pm + p_{dc})e^{-(\lambda_\pm+p_{dc})}
% 	\label{eq:prob_click}
% \end{equation}
%
% \subsection{2 Detector Statistics}
% As mentioned before, we have 2 detectors with the same efficiency $\eta$ and the same detector dark rate $p_{dc}$. These can produce a click independent from one another, but we're only interested in the cases where \textbf{ONLY} 1 produces a click, as the photon can (ideally) only be either $|+\rangle$ or $|-\rangle$ polarized.
% \begin{align*}
% 	P(\text{only}|+\rangle) & = P_{+,click}\cdot P_{-,0 clicks}                                                \\
% 	                        & = (\lambda_{+} + p_{dc})e^{-(\eta(\bar{n}_{out,+} + \bar{n}_{out,-}) + 2p_{dc})}
% \end{align*}
% \begin{align*}
% 	P(\text{only}|-\rangle) = (\lambda_{-} + p_{dc})e^{-(\eta(\bar{n}_{out,+} + \bar{n}_{out,-}) + 2p_{dc})}
% \end{align*}
%
% Adding the two together gives us the probability that at just from the fact that only one detector clicked, something sort of useful came out of our output field, which we're going to call $P_{valid}$:
% $$
% 	P_{valid} = (\lambda_{+} + \lambda_{-} + 2p_{dc})e^{-(\eta(\bar{n}_{out,+}+\bar{n}_{out,-})+2p_{dc})}
% $$
%
% However now we need to filter out only the cases where the detector event was actually triggered by a photon and not a dark count.
%
% \subsection{Photon Event Filtering}
% Here the probability that the 1 detector event in \textbf{ONLY} was caused by a photon is:
% \begin{align*}
% 	P(\text{1 photon in + \& no click in -}) & = \lambda_{+} e^{-(\lambda_{+} + p_{dc})} \cdot e^{-(\lambda_{-} + p_{dc})} \\
% 	                                         & = \lambda_{+} e^{-(\eta(\bar{n}_{out,+}+\bar{n}_{out,-})+2p_{dc})}          \\
% 	P(\text{1 photon in - \& no click in +}) & =  \lambda_{-} e^{-(\eta(\bar{n}_{out,+}+\bar{n}_{out,-})+2p_{dc})}
% \end{align*}
%
% Thus giving us the probability that 1 click was produced in only 1 of the two detectors:
% $$ P_{only} = (\lambda_{+} + \lambda_{-}) e^{-(\eta(\bar{n}_{out,+}+\bar{n}_{out,-})+2p_{dc})}$$
%
% This then allows us now to calculate how many of our valid clicks, were actually produced by only 1 photon event going only to 1 of the 2 detectors, essentially giving us a sort of signal-to-noise ratio:
% $$
% 	P_{SNR} = \frac{P_{only}}{P_{valid}} = \frac{\lambda_{+} + \lambda_{-}}{\lambda_{+} + \lambda_{-} + 2p_{dc}}
% $$
%
% \subsection{Photon Distribution}
% All that's left now is to calculate how much of our light went to the "+"-detector, and how much went to our "-"-detector, which is just like calculating the ratios of all of our events that were triggered only by photons, and the probability of the respective detector itself:
% $$
% 	R_{+} = \frac{P(\text{1 photon in + and no click in -})}{P_{only}} = \frac{\bar{n}_{out,+} }{\bar{n}_{out,+} + \bar{n}_{out,-}}$$
% $$
% 	R_{-} = \frac{P(\text{1 photon in - and no click in +})}{P_{only}} = \frac{\bar{n}_{out,-}}{\bar{n}_{out,+} + \bar{n}_{out,-}}
% $$
%
% What then ends up going into our truth table is the following:
% \begin{itemize}
% 	\item For the $|+\rangle$-part = $R_{+}\cdot P_{SNR} + (1-P_{SNR})\cdot 0.5$
% 	\item For the $|-\rangle$-part = $R_{-}\cdot P_{SNR} + (1-P_{SNR})\cdot 0.5$
% \end{itemize}
%
% \section{Driving Hamiltonian}
%
% Since we're mainly going to be working with the \textit{mesolve} function in QuTip to simulate the dynamics of the atom-photon gate it's imperative to know what each and every part does. Next to the Jaynes-Cummings-Hamiltonian the most important part is the one explaining the \textit{photon that gets reflected from the cavity}, most notably the part before it gets reflected aka the \textit{driving Hamiltonian} $H_{d}$. For this we're going to just write down once, the dynamics of a coherent driving field which will be the foundation of my photon that shall or shall not enter the cavity.
% \newline
%
% The photon is going to come from a laser, that is going to be described as a classical oscillating electric field. That light field has an amplitude $E_0$ and a frequency $\omega_d$:
% $$
% 	E(t) = E_0 \cos(\omega_d t)
% $$
% Using the euler formula we can rewrite this to:
% $$
% 	E(t) = \frac{E_0}{2} (e^{i\omega_d t} + e^{-i\omega_d t})
% $$
%
% The light field operator $\hat{E}$ interacts with the light in our cavity (which get described by our creation and annihilation operators) linearly:
% $$
% 	\hat{E} \propto (a + a^\dag)
% $$
%
% When now describing the \textit{driving Hamiltonian} what we're doing is just overlaping the light field operator with the oscillating electric field:
% \begin{align*}
% 	H_{drive} & \propto E(t) \cdot (a+a^{\dag})                         \\
% 	H_{drive} & = (e^{i\omega_d t} + e^{-i\omega_d t})\cdot(a+a^{\dag})
% \end{align*}
%
% From this we get four terms where two cancel out as they aren't energy conserving rendering us with the final expression for our driving Hamiltonian:
% \begin{empheq}[box=\mymath]{equation*}
% 	H_{drive} = a^\dag e^{-i\omega_d t} + a e^{+i\omega_d t}
% \end{empheq}
%
% Where the first term adds a photon to our cavity and the other removes one. Sometimes you'll see a minus instead of a plus between the two terms, this is just due to different conventions of the light field, where people chose a $\sin$ to describe the oscillating part instead of a $\cos$, but the physics remain the same.
%
% \section{Analytical derivation for case of atom in $\ket{0}$ and photon is in $\ket{h}$}
%
% Our Hamiltonian is as follows:
% \begin{align*}
% 	H = \hbar g(\ket{e}\bra{1}a_h + \ket{1}\bra{e}a_h^{\dagger})
% \end{align*}
%
% We drive the $a_h$ mode of our cavity through $a_h^{in}(t)$, thus giving us the following equation of motion for our $a_h$ cavity mode:
%
% \begin{align*}
% 	\dot{a_h} & =-i[a_h,H]-(i\Delta + \frac{\kappa}{2})a_h - \sqrt{\kappa}a_h^{in}(t) \\
% \end{align*}
% Because $a_h$ commutes with $H$, the commutator is 0 and we get
% \begin{align*}
% 	\rightarrow \dot{a_h} & =-(i\Delta + \frac{\kappa}{2})a_h - \sqrt{\kappa}a_h^{in}(t)
% \end{align*}
%
% We now drop the subscript $h$ as we'll only be looking at that case.
%
% Multiplying with $e^{i\omega t}$, and integrating we get:
% \begin{align*}
% 	\int{dt e^{i\omega t}\dot{a}(t)} =  \underbrace{[e^{i\omega t}a(t)]_{-\infty}^{\infty}}_{=0} - \underbrace{i\omega\int{dt e^{i\omega t}a(t)}}_{=i\omega a(\omega)}
% \end{align*}
% where in the last expression we used that the integral performs a Fourier transformation.
%
% \begin{align*}
% 	\Rightarrow -i\omega a(\omega) & = -(i\Delta + \frac{\kappa}{2})a(\omega) - \sqrt{\kappa}a^{in}(\omega)  \\
% 	\sqrt{\kappa}a^{in}(\omega)    & = -(i\Delta + \frac{\kappa}{2} - i\omega)a(\omega)                      \\
% 	\rightarrow a(\omega)          & = -\frac{\sqrt{\kappa}}{i\Delta+\frac{\kappa}{2}-i\omega}a^{in}(\omega)
% \end{align*}
%
% Now from our regular input output relation
% \begin{align*}
% 	a^{out} (\omega)= a^{in}(\omega) + \kappa a(\omega)
% \end{align*}
%
% We thus get:
%
% \begin{align*}
% 	a^{out} (\omega) & =a^{in}(\omega) -\frac{\sqrt{\kappa}}{i\Delta+\frac{\kappa}{2}-i\omega}a^{in}(\omega)                               \\
% 	a^{out}(\omega)  & =(1-\frac{\kappa}{i\Delta+\frac{\kappa}{2}-i\omega})a^{in}(\omega)                                                  \\
% 	a^{out}(\omega)  & =(1-\frac{i\Delta + \frac{\kappa}{2}-i\omega-\kappa}{i\Delta+\frac{\kappa}{2}-i\omega})a^{in}(\omega)               \\
% 	a^{out}(\omega)  & =(\frac{i\Delta-\frac{\kappa}{2}-i\omega}{i\Delta+\frac{\kappa}{2}-i\omega})a^{in}(\omega)                          \\
% 	                 & =\underbrace{\frac{i(\Delta-\omega)-\frac{\kappa}{2}}{i(\Delta-\omega)+\frac{\kappa}{2}}}_{h(\omega)}a^{in}(\omega)
% \end{align*}
%
% We can then use the \href{https://en.wikipedia.org/wiki/Convolution_theorem}{convolution theorem} in order to solve for $a^{out}(t)$
%
% \begin{align*}
% 	a^{out}(t) = \int_{-\infty}^{\infty}d\tau \underbrace{h(t-\tau)}_{h(t) =\frac{1}{2\pi}\int_{-\infty}^{\infty}d\omega e^{-i\omega t}h(\omega)}a^{in}(\tau)
% \end{align*}
%
% Now using that $h(\omega) = 1-\frac{\kappa}{i(\Delta-\omega)+\frac{\kappa}{2}}$ we get for $h(t)$ (via partial integration):
%
% \begin{align*}
% 	h(t)                   & = \delta(t) - \kappa e^{-i\delta t}e^{-\frac{\kappa}{2}t}\theta(t)                                          \\
% 	\Rightarrow a^{out}(t) & = a^{in}(t) - \kappa\int_{-\infty}^{\infty}d\tau e^{-i\delta(t-\tau)}e^{-\frac{\kappa}{2}(t-\tau)}a^{in}(t)
% \end{align*}
%
% For slowly varying pulses we can now approximate that $a^{in}(\tau)\approx a^{in}(t)$ thus giving us:
%
% \begin{align*}
% 	\Rightarrow a^{out}(t) & \approx a^{in}(t) - \frac{\kappa}{i\Delta+\frac{\kappa}{2}}a^{in}(t) \\
% 	                       & =\frac{i\Delta-\frac{\kappa}{2}}{i\Delta+\frac{\kappa}{2}}a^{in}(t)
% \end{align*}
%
% Which for no detuning (so $\Delta=0$) gives us:
% \begin{equation}
% 	a^{out}(t) =-a^{in}(t)
% \end{equation}
%
% \subsection{Lindblad Formalism}
%
% We start from the master equation:
% \begin{equation}
% 	\dot{\rho} = -\frac{i}{\hbar}[H,\rho] + \sum_j \gamma_j \left( L_j \rho L_j^\dagger - \tfrac{1}{2}\{ L_j^\dagger L_j, \rho \} \right),
% \end{equation}
%
% Now, for an operator \(O\),
% \begin{align}
% 	\frac{d}{dt} \langle O \rangle & = \Tr \left( O \dot{\rho} \right)                                                                                                                               \\
% 	                               & = \Tr \left( O \left[ -\frac{i}{\hbar}[H,\rho] + \sum_j \gamma_j \left( L_j \rho L_j^\dagger - \tfrac{1}{2}\{ L_j^\dagger L_j, \rho \} \right) \right] \right).
% \end{align}
%
% The Hamiltonian contribution gives
% \begin{equation}
% 	- \Tr \left( O \frac{i}{\hbar}[H,\rho] \right)
% 	= -\frac{i}{\hbar}\Tr \left( O H \rho - O \rho H \right)
% 	= -\frac{i}{\hbar}\Tr \left( [O,H]\rho \right)
% 	= -\frac{i}{\hbar} \langle [O,H] \rangle.
% \end{equation}
%
% For the dissipative terms we obtain
% \begin{align}
% 	\sum_j \gamma_j \Tr\Big( O L_j \rho L_j^\dagger
% 	- \tfrac{1}{2} O L_j^\dagger L_j \rho
% 	- \tfrac{1}{2} O \rho L_j^\dagger L_j \Big).
% \end{align}
%
% Using cyclicity of the trace:
% \begin{align}
% 	 & = \sum_j \gamma_j \left( \langle L_j^\dagger O L_j \rangle
% 	- \tfrac{1}{2} \langle \{ L_j^\dagger L_j, O \} \rangle \right).
% \end{align}
%
% Therefore the general equation of motion is
% \begin{equation}
% 	\frac{d}{dt}\langle O \rangle
% 	= -\frac{i}{\hbar}\langle [O,H] \rangle
% 	+ \sum_j \gamma_j \left( \langle L_j^\dagger O L_j \rangle
% 	- \tfrac{1}{2}\langle \{L_j^\dagger L_j, O \} \rangle \right).
% \end{equation}
%
%
% \subsection{Calculation of the Maxwell Bloch Equations (no lossy channels)}
% \subsubsection{Jaynes Cummings Hamiltonian}
% Assuming ground state has zero energy our Jaynes-Cummings Hamiltonian is equal to:
% \begin{equation}
% 	H_{JC,S} = \underbrace{\omega_ca^\dagger a + \omega_a\sigma^\dagger\sigma}_{=H_0,S} + \underbrace{g(\sigma + \sigma^\dagger)(a +  a^\dagger)}_{=H_{int,S}} + H_{drive,S}
% \end{equation}
% where $\sigma$ is the atomic lowering operator,$\omega_a$ the atomic transition frequency between the levels $\ket{1}$ and $\ket{e}$,$a$ the photonic destruction operator, and $\omega_a$ the cavity mode resonance frequency.\\
% $H_{drive}$ in this case is the external light field which is driving our system, which in the semiclassical formalism has the following form \cite{WallsMilburn2008}:
% \begin{equation}
% 	H_{drive,S} = i(\epsilon e^{-i\omega_d t} + \epsilon^* e^{i\omega_d t})(a + a^\dagger)
% \end{equation}
%
% Here, $a$ again is the photonic destruction operator, $\omega_d$ is the frequency of the driving field, and $\epsilon$ is the complex amplitude of our incoming pulse.
%
% In general going from the Schrödinger picture to the interaction picture, is achieved by converting our operators into the interaction picture via following equation:
% \begin{equation}
% 	i\frac{dA_I}{dt} = [A_I(t),H_{0,S}]
% \end{equation}
%
% Furthermore, $H_{0,I} = H_{0,S}$. Equipped with this we can now calculate the expressions for operators in the interaction picture:
%
% \begin{align}
% 	i\frac{da_I(t)}{dt}                                & = [a_I(t), H_{0,S}]                    \\
% 	                                                   & = [a_I(t),\omega_ca^\dagger a]         \\
% 	                                                   & = \omega_c a_I(t)                      \\
% 	\Rightarrow \frac{da_I(t)}{dt} = -i\omega_c a_I(t) & \rightarrow a_I(t) = a e^{-i\omega_ct}
% \end{align}
%
% Where we used that:
% \begin{equation}
% 	[a,H_{0,S}] = \omega_c[a,a^\dagger a] = \omega_c (\underbrace{[a,a^\dagger]}_{=1}a + a^\dagger\underbrace{[a,a]}_{=0}) = \omega_c a
% \end{equation}
%
% Now we repeat the same thing for the atomic lowering operator $\sigma$:
%
% \begin{align}
% 	i\frac{d\sigma_I(t)}{dt}                                    & = \underbrace{[\sigma_I(t),H_{0,S}]}             \\
% 	i\frac{d\sigma_I(t)}{dt}                                    & = \omega_a\sigma_I(t)                            \\
% 	\Rightarrow \frac{d\sigma_I(t)}{dt} = -i\omega_a\sigma_I(t) & \rightarrow \sigma_I(t) = \sigma e^{-i\omega_at}
% \end{align}
%
% Now we can calculate the Hamiltonian for the interaction part and the drive part, in the interaction picture, by just plugging these operators in there, giving us:
%
% \begin{align}
% 	H_{int,I}         & =\sigma a e^{-i(\omega_a-\omega_c)t} + \sigma a^\dagger e^{-i(\omega_a-\omega_c)t} +\sigma^\dagger a e^{-i(\omega_c-\omega_a)t} + \sigma^\dagger a^\dagger e^{i(\omega_c+\omega_a)t} \\
% 	\xrightarrow{RWA} & = \sigma a^\dagger e^{-i(\omega_a-\omega_c)t} +\sigma^\dagger a e^{-i(\omega_c-\omega_a)t}                                                                                           \\
% 	H_{drive,I}       & = i(\epsilon e^{-i\omega_d t} + \epsilon^\dagger e^{i\omega_d t})(a_I(t) + a_I^\dagger(t))                                                                                           \\
% 	\xrightarrow{RWA} & = i(a \epsilon^*e^{-i(\omega_c-\omega_d)t} + a^\dagger \epsilon e^{i(\omega_c-\omega_d)t})
% \end{align}
%
% Thus giving us the total Hamiltonian in the interaction picture under the rotating wave approximation:
%
% \begin{equation}
% 	H_{JC,I}\overset{\text{RWA}}{=} \omega_ca^\dagger a + \omega_a \sigma^\dagger\sigma + g(\sigma a^\dagger e^{-i(\omega_a-\omega_c)t} + \sigma^\dagger a e^{-i(\omega_c-\omega_a)t}) + i(a \epsilon^* e^{-i(\omega_c-\omega_d)t} + a^\dagger \epsilon e^{i(\omega_c-\omega_d)t})
% \end{equation}
%
% If we now transform back into the Schrödinger picture, we get:
% \begin{equation}
% 	H_{JC,S} \overset{\text{RWA}}{=} \omega_ca^\dagger a + \omega_a \sigma^\dagger\sigma + g(a^\dagger\sigma + a\sigma^\dagger) + i(a \epsilon^*e^{i\omega_dt} + a^\dagger \epsilon e^{-i\omega_dt})
% \end{equation}
%
% If we now change into the rotating frame of the driving field via
% \begin{align}
% 	U(t) & = \exp{i\omega_dt(a^\dagger a + \sigma^\dagger\sigma)} \\
% 	H    & \rightarrow UHU^\dagger + i\dot{U}U^\dagger =: H_{RF}
% \end{align}
%
% we then get the final form of our Hamiltonian:
% \begin{empheq}[box=\fcolorbox{red}{yellow!10}]{equation}
% 	H_{JC,RF} = \underbrace{(\omega_c-\omega_d)}_{=\Delta_c}a^\dagger a
% 	+ \underbrace{(\omega_a-\omega_d)}_{=\Delta_a}\sigma^\dagger\sigma
% 	+ g(a^\dagger\sigma + a\sigma^\dagger)
% 	+ i(a\epsilon^* + a^\dagger\epsilon)
% \end{empheq}
%
% Which gives rise to 4 scenarios which I'll be covering with my project:
% \begin{center}
% 	\resizebox{\linewidth}{!}{
% 		\begin{tabular}{|l|l|}
% 			\hline
% 			\textbf{scenario}          & \textbf{meaning}                                                                                                 \\
% 			\hline
% 			$\Delta_a=\Delta_c=0$      & \makecell{Atom is in coupled state (here: $\ket{1}$ and we drive with a light field                              \\ that's resonant to the cavity (here: we drive $\ket{1} \leftrightarrow \ket{e}$ which would be our $\ket{\pi}$ polarized light}  \\
% 			\hline
% 			$\Delta_a=0;\Delta_c\neq0$ & Atom is in coupled state but we drive with light that's non-resonant to cavity (here: $\ket{V}$ polarized light) \\
% 			\hline
% 			$\Delta_a\neq0;\Delta_c=0$ & Atom is in uncoupled state (here: $\ket{0}$) and we drive with a light field that's resonant to the cavity       \\
% 			\hline
% 			$\Delta_a,\Delta_c\neq0$   & Atom is in uncoupled state and we drive with light that's non-resonant to cavity                                 \\
% 			\hline
% 		\end{tabular}}
% \end{center}
%
% In order to improve readability we'll be splitting the full JC Hamiltonian into three parts:
% \begin{align}
% 	H_0       & = \Delta_ca^\dagger a + \Delta_a\sigma^\dagger\sigma \\
% 	H_{int}   & = g(a^\dagger\sigma + a\sigma^\dagger)               \\
% 	H_{drive} & = i(a\epsilon^* + a^\dagger\epsilon)
% \end{align}
%
% with following losses:
% \begin{align}
% 	\gamma_1 & = \gamma, \quad L_1 = \sigma, \\
% 	\gamma_2 & = \kappa, \quad L_2 = a.
% \end{align}
%
%
% \subsection{Cavity mode a}
%
% For $H_0$:
% \begin{align}
% 	[a,H_0] & = \Delta_c[a,a^\dagger a] + \Delta_ca \underbrace{[a,\sigma^\dagger\sigma]}_{=0} \\
% 	        & = \Delta_c (\underbrace{[a,a^\dagger]}_{=1}a + a^\dagger\underbrace{[a,a]}_{=0}) \\
% 	        & = \Delta_c a
% \end{align}
% \begin{equation}
% 	\boxed{
% 		\Rightarrow -i\langle[a,H_0]\rangle = -i\Delta_c \langle a\rangle
% 	}
% \end{equation}
% For $H_{int}$:
% \begin{align}
% 	[a,H_{int}] & = g([a,a^\dagger\sigma] + [a,a\sigma^\dagger])                                                           \\
% 	            & = g(a a^\dagger\sigma - a^\dagger a\sigma + \cancel{a a\sigma^\dagger} - \cancel{aa\sigma^\dagger})      \\
% 	            & \rightarrow \langle aa^\dagger\sigma\rangle = \bra{n}aa^\dagger\sigma\ket{n} =(n+1) \langle\sigma\rangle \\
% 	            & \rightarrow \langle a^\dagger a\sigma\rangle = \bra{n}a^\dagger a\sigma\ket{n} =n \langle\sigma\rangle
% \end{align}
% \begin{equation}
% 	\boxed{
% 	\Rightarrow -i\langle[a,H_{int}]\rangle = -ig\langle\sigma\rangle
% 	}
% \end{equation}
% For $H_{drive}$:
% \begin{align}
% 	[a,H_{drive}] & = [a,i(a\epsilon^* + a^\dagger\epsilon)]     \\
% 	              & = i([a,a\epsilon^*] + [a,a^\dagger\epsilon]) \\
% 	              & = i(0 + \epsilon)                            \\
% 	              & = i\epsilon
% \end{align}
% \begin{equation}
% 	\boxed{
% 	\Rightarrow -i\langle[a,H_{drive}]\rangle = i\langle \epsilon\rangle \overset{\epsilon \in \mathds{C}}{=} i\epsilon
% 	}
% \end{equation}
% If we now include the fact that:
% \begin{equation}
% 	i\epsilon \mathrel{\widehat{=}} \sqrt{\kappa_{oc}}a^{in}(t)
% \end{equation}
% we are then left with:
% \begin{equation}
% 	\boxed{
% 	\Rightarrow-i\langle[a,H_{drive}]\rangle = -\sqrt{\kappa_{oc}}a^{in}(t)
% 	}
% \end{equation}
%
% For rest of the Lindbladian:
% \begin{align}
% 	D(a) & =a^\dagger a a - \frac{1}{2}\{a^\dagger a,a\}               \\
% 	     & =a^\dagger a a - \frac{1}{2}(a^\dagger a a - a a^\dagger a) \\
% 	     & = \frac{1}{2}a^\dagger a a - \frac{1}{2} a a^\dagger a      \\
% 	     & = \frac{1}{2}(a^\dagger a a - a a^\dagger a)                \\
% 	     & = \frac{1}{2}([a^\dagger,a]a)                               \\
% 	     & = \frac{1}{2}(-a)                                           \\
% 	     & = -\frac{1}{2} a
% \end{align}
% \begin{equation}
% 	\boxed{
% 		\Rightarrow \kappa\langle D(a)\rangle = -\frac{\kappa}{2} \langle a\rangle
% 	}
% \end{equation}
%
% So for $\langle a\rangle(t)$ we get:
% \begin{empheq}[box=\fcolorbox{red}{yellow!10}]{align}
% 	\frac{d\langle a\rangle}{dt}
% 	&= -i\Delta_c\langle a\rangle - i g \langle\sigma\rangle
% 	- \sqrt{\kappa_{oc}}\, a_{\text{in}}(t) - \frac{\kappa}{2}\,\langle a\rangle \notag\\
% 	&= -(i\Delta_c + \tfrac{\kappa}{2})\langle a\rangle
% 	- i g \langle\sigma\rangle - \sqrt{\kappa_{oc}}\, a_{\text{in}}(t)
% \end{empheq}
%
% \subsection{Atomic lowering operator $\sigma$}
% For $H_0$:
% \begin{align}
% 	[\sigma,H_0] & = \Delta_a\underbrace{[\sigma,\sigma^\dagger\sigma]}_{=\underbrace{[\sigma,\sigma^\dagger]}_{=\mathds{1}}\sigma + \sigma^\dagger\underbrace{[\sigma,\sigma]}_{=0}} \\
% 	             & = \Delta_a \mathds{1}\sigma
% \end{align}
% \begin{equation}
% 	\boxed{
% 		\Rightarrow -i\langle[\sigma,H_0]\rangle = -i\Delta_a \langle \sigma\rangle
% 	}
% \end{equation}
% For $H_{int}$:
% \begin{align}
% 	[\sigma,H_{int}] & = g([\sigma,a^\dagger\sigma] + [\sigma,a\sigma^\dagger])                                                              \\
% 	                 & = g(\cancel{a^\dagger \sigma\sigma} - \cancel{a^\dagger\sigma\sigma} + a\sigma\sigma^\dagger - a\sigma^\dagger\sigma) \\
% 	                 & = g(a\sigma\sigma^\dagger - a\sigma^\dagger\sigma)
% \end{align}
% Now we can consider that:
% \begin{empheq}[box=\fcolorbox{red}{yellow!10}]{align}
% 	\sigma^\dagger\sigma &= \ket{e}\braket{1}\bra{e} = \ket{e}\bra{e} \\
% 	\sigma\sigma^\dagger &= \ket{1}\braket{e}\bra{1} = \ket{1}\bra{1}
% \end{empheq}
% \begin{align}
% 	[\sigma,H_{int}] & = g(a\sigma\sigma^\dagger - a\sigma^\dagger\sigma) \\
% 	                 & = ga(\ket{1}\bra{1} - \ket{e}\bra{e})              \\
% 	                 & = ga \sigma_z
% \end{align}
% \begin{equation}
% 	\boxed{\Rightarrow -i\langle[\sigma,H_{int}]\rangle=-ig\langle a\sigma_z \rangle}
% \end{equation}
% No interaction with $H_{drive}$.\\
% For rest of Lindbladian:
% \begin{align}
% 	D(\sigma) & = \underbrace{\sigma^{\dagger}\underbrace{\sigma\sigma}_{=0}}_{=0} - \frac{1}{2} \{\sigma^\dagger\sigma,\sigma\} \\
% 	          & = -\frac{1}{2}\underbrace{\sigma^\dagger\sigma\sigma}_{=0} - \frac{1}{2}\sigma\sigma^\dagger\sigma               \\
% 	          & = -\frac{1}{2}\sigma
% \end{align}
% \begin{equation}
% 	\boxed{\gamma\langle D(\sigma)\rangle = -\frac{\gamma}{2}\langle\sigma\rangle}
% \end{equation}
%
% So for $\langle\sigma\rangle$(t) we get:
% \begin{empheq}[box=\fcolorbox{red}{yellow!10}]{align}
% 	\frac{d\langle\sigma\rangle}{dt} &= -i\Delta_a\langle\sigma\rangle -ig\langle a\sigma_z\rangle -\frac{\gamma}{2}\langle\sigma\rangle \notag\\
% 	&= -(i\Delta_a + \frac{\gamma}{2})\langle\sigma\rangle - ig\langle a\sigma_z\rangle
% \end{empheq}
%
% \subsection{Pauli-Z operator $\sigma_z$}
%
% For $H_0$:
% \begin{align}
% 	[\sigma_z,H_0] & = \Delta_a[\sigma_z,\sigma^\dagger\sigma]                     \\
% 	               & = \Delta_a[\sigma_z,\ket{e}\bra{e}]                           \\
% 	               & = \Delta_a (\sigma_z \ket{e}\bra{e} - \ket{e}\bra{e}\sigma_z) \\
% 	               & = \Delta_a (-\ket{e}\bra{e} + \ket{e}\bra{e})                 \\
% 	               & = 0
% \end{align}
% For $H_{int}$:
% \begin{align}
% 	[\sigma_z, H_{int}] & = g([\sigma_z,a^\dagger\sigma] + [\sigma_z,a\sigma^\dagger])                                               \\
% 	                    & = g(a^\dagger\sigma_z\sigma - a^\dagger\sigma\sigma_z + a\sigma_z\sigma^\dagger - a\sigma^\dagger\sigma_z) \\
% 	                    & = g(a^\dagger\sigma + a^\dagger\sigma -a\sigma^\dagger - a\sigma^\dagger)                                  \\
% 	                    & = 2g(a^\dagger\sigma - a\sigma^\dagger)
% \end{align}
% \begin{equation}
% 	\boxed{
% 	\Rightarrow -i\langle[\sigma_z,H_{int}]\rangle=-2ig(\langle a^\dagger\sigma \rangle - \langle a\sigma^\dagger\rangle)
% 	}
% \end{equation}
% No interaction with $H_{drive}$.
% For rest of Lindbladian:
% \begin{align}
% 	D(\sigma_z) & = \sigma^\dagger\sigma_z\sigma - \frac{1}{2}\{\sigma^\dagger\sigma,\sigma\}                                                                                           \\
% 	            & = \ket{e}\bra{e}  - \frac{1}{2}\underbrace{\sigma^\dagger\sigma\sigma_z}_{=-\ket{e}\bra{e}} - \frac{1}{2}\underbrace{\sigma_z\sigma^\dagger\sigma}_{=-\ket{e}\bra{e}} \\
% 	            & = \ket{e}\bra{e} + \ket{e}\bra{e}                                                                                                                                     \\
% 	            & = 2\ket{e}\bra{e}                                                                                                                                                     \\
% 	            & = \mathds{1} - \sigma_z
% \end{align}
% \begin{equation}
% 	\boxed{
% 		\Rightarrow\gamma\langle D(\sigma_z)\rangle=\gamma(1-\langle\sigma_z\rangle)
% 	}
% \end{equation}
%
% So for $\langle\sigma_z\rangle(t)$ we get:
% \begin{empheq}[box=\fcolorbox{red}{yellow!10}]{equation}
% 	\frac{d\langle\sigma_z\rangle(t)}{dt} = -2ig(\langle a^\dagger\sigma \rangle - \langle a\sigma^\dagger\rangle) + \gamma(1-\langle\sigma_z\rangle)
% \end{empheq}
%
% The final Maxwell Bloch Equations then take the form of:
% \begin{empheq}[box=\fcolorbox{red}{yellow!10}]{align}
% 	\frac{d\langle a\rangle}{dt} &= -(i\Delta_c + \tfrac{\kappa}{2})\langle a\rangle
% 	- i g \langle\sigma\rangle - \sqrt{\kappa_{oc}}\, a_{\text{in}}(t) \\
% 	\frac{d\langle\sigma\rangle}{dt} &= -(i\Delta_a + \tfrac{\gamma}{2})\langle\sigma\rangle
% 	- i g \langle a\sigma_z\rangle \\
% 	\frac{d\langle\sigma_z\rangle}{dt} &= -2i g \big(\langle a^\dagger\sigma \rangle - \langle a\sigma^\dagger\rangle\big)
% 	+ \gamma \big(1 - \langle\sigma_z\rangle\big)
% \end{empheq}
